\chapter{Dark Circles Under Eyes}

\section*{Definition}
Dark, discolored appearance of the skin beneath the eyes, which may be black, blue, brown, or purple, resulting from vascular, pigmentary, or structural causes.

\section*{What Happens in Your Body}
Dark circles arise from multiple mechanisms: vascular congestion (thin periorbital skin reveals underlying blood vessels, exacerbated by fatigue, allergies, or aging), melanin deposition (constitutional/genetic hyperpigmentation, more common in darker skin types), post-inflammatory hyperpigmentation (from atopic or allergic dermatitis), and shadowing from infraorbital fat herniation or tear trough hollowing with aging.

\section*{Causes}
\begin{itemize}
  \item Fatigue and sleep deprivation (vascular stasis)
  \item Allergic rhinitis and atopic dermatitis (venous congestion)
  \item Genetic/constitutional periorbital hyperpigmentation
  \item Aging (skin thinning, fat herniation, tear trough hollowing)
  \item Dehydration
  \item Iron deficiency anemia
  \item Sun exposure (increased melanin production)
  \item Eye strain or prolonged screen time
  \item Periorbital edema (fluid retention)
  \item Medications (chemical messengers that cause inflammation and pain analogs for glaucoma)
\end{itemize}

\section*{When to See a Doctor}
See your doctor if:
\begin{itemize}
  \item Symptoms are severe or getting worse over time
  \item Dark circles with persistent fatigue, pallor, or other signs of anemia, or if accompanied by periorbital swelling or pain
  \item You have other concerning symptoms such as fever, weight loss, or bleeding
\end{itemize}

\section*{Related Symptoms}
\begin{itemize}
  \item Eye puffiness
  \item Fatigue
  \item Nasal congestion
  \item Headache
  \item Pale skin
\end{itemize}

\textbf{Important:} If this symptom persists, your doctor will check for serious underlying causes.

\section*{Self-Care / Home Management}
\begin{itemize}
  \item Ensure adequate sleep (7-9 hours) and maintain a consistent sleep schedule.
  \item Apply cold compresses or chilled cucumber slices to reduce periorbital puffiness.
  \item Use topical creams containing vitamin C, retinol, or kojic acid to address pigmentation.
  \item Elevate your head while sleeping to reduce fluid accumulation under the eyes.
  \item Manage allergies with antihistamines and avoid known triggers.
\end{itemize}

\section*{How Your Doctor Will Evaluate You}
\begin{itemize}
  \item Your doctor will examine the periorbital area and assess skin thickness, pigmentation, and vascular prominence.
  \item A review of your sleep habits, allergy history, and skincare routine will be conducted.
  \item Blood tests may check for iron deficiency anemia, thyroid function, and vitamin levels.
  \item Allergy testing may be recommended if allergic shiners are suspected.
  \item Referral to a dermatologist or allergist may be arranged for persistent or severe cases.
\end{itemize}

\section*{Treatment Options}
\begin{itemize}
  \item Addressing underlying causes such as treating allergies, correcting iron deficiency, or improving sleep quality.
  \item Topical depigmenting agents (hydroquinone, azelaic acid) for melanin-based dark circles.
  \item Chemical peels or laser therapy for resistant pigmentation, performed by a dermatologist.
  \item Dermal fillers for tear trough hollowing causing shadow-related dark circles.
  \item A consistent skincare routine with sunscreen, moisturizer, and targeted serums.
\end{itemize}

\section*{References}
\begin{thebibliography}{99}
\bibitem{ref1} Sarkar R, Sinha S. "Periorbital hyperpigmentation: a comprehensive review." J Cutan Aesthet Surg. 2015;8(4):197-205.
\bibitem{ref2} Ranu H, Thng S, Goh BK, et al. "Periorbital hyperpigmentation in Asians: a review of current understanding and treatment." Dermatol Surg. 2011;37(10):1479-1488.
\bibitem{ref3} Gupta M, Prakash S, Mahajan VK. "Periorbital hyperpigmentation: etiopathogenesis and treatment." Indian Dermatol Online J. 2020;11(5):697-705.
\bibitem{ref4} Sheth PB, Shah HA, Dave JN. "Periorbital hyperpigmentation: a study of its prevalence, common causative factors and its association with personal habits and systemic diseases." Indian J Dermatol. 2014;59(2):150-155.
\bibitem{ref5} Alexis AF, Alam M. "Cosmetic dermatology for skin of color." Dermatol Clin. 2023;41(3):457-468.
\bibitem{ref6} Grimes PE. "Periorbital hyperpigmentation." UpToDate. Wolters Kluwer; 2025.
\end{thebibliography}
